% 2.2.2
% Ficheiro incluído por main.tex

\subsection[Preparar o cliente: desativar o SSSD e ativar o Winbind]{Preparar o cliente: desativar o SSSD e ativar o Winbind}

\esquema{2-2-2-trocar-variante}{Neste cliente desativa-se o SSSD e ativa-se o Winbind}

\textbf{Objetivo:} este cliente foi clonado do mesmo template UbuntuClient, que vem com o SSSD ativo e o Winbind desativado. Aqui faz-se o contrário: desativa-se o SSSD e ativa-se o Winbind.

\passo{Passo 1 - Parar e desativar o serviço SSSD}

\begin{lstlisting}[style=bashstyle]
sudo systemctl disable --now sssd
systemctl is-enabled sssd
\end{lstlisting}

\emph{Verificação esperada:} \texttt{disabled}.

\passo{Passo 2 - Retirar o SSSD do PAM}

\begin{lstlisting}[style=bashstyle]
sudo pam-auth-update --disable sss
grep -c pam_sss /etc/pam.d/common-auth
\end{lstlisting}

\emph{Verificação esperada:} \texttt{0} (o \texttt{pam\_sss} já não está na pilha de autenticação).

\passo{Passo 3 - Retirar o SSSD do NSS}

\begin{lstlisting}[style=bashstyle]
sudo cp -a /etc/nsswitch.conf /etc/nsswitch.conf.orig
sudo sed -i -E '/^(passwd|group|shadow):/ s/ sss//' /etc/nsswitch.conf
grep -E '^(passwd|group|shadow):' /etc/nsswitch.conf
\end{lstlisting}

\emph{Verificação esperada:} as linhas \texttt{passwd:}, \texttt{group:} e \texttt{shadow:} já não têm \texttt{sss}. O \texttt{winbind} é configurado no NSS em 2.2.7.

\passo{Passo 4 - Ativar o serviço Winbind}

\begin{lstlisting}[style=bashstyle]
sudo systemctl enable winbind
systemctl is-enabled winbind
\end{lstlisting}

\emph{Verificação esperada:} \texttt{enabled}. O serviço só arranca depois de configurado o \texttt{smb.conf} e feito o join (2.2.5, Passo 1).

\problemas{\hyperref[ec-b2]{Erros Comuns - 2.2 Winbind}}
